\documentclass[UTF8,zihao=-4]{ctexart}
\usepackage[a4paper,margin=1.85cm]{geometry}
\usepackage{hyperref}
\usepackage{verbatim}
\usepackage{xcolor}
\usepackage{fancyhdr}
\hypersetup{hidelinks}
\setlength{\parindent}{0pt}
\setlength{\parskip}{0.45em}
\setlength{\headheight}{14.5pt}
\emergencystretch=4em
\sloppy
\pagestyle{fancy}
\fancyhf{}
\lhead{编译原理专项模拟卷}
\rhead{\thepage}
\title{覆盖矩阵与组卷说明}
\author{}
\date{}
\begin{document}
\maketitle
\setcounter{tocdepth}{1}
\tableofcontents
\newpage

\section*{一、使用顺序}
\addcontentsline{toc}{section}{一、使用顺序}


本目录分成两组使用：

\subsection*{可独立完成的完整试卷}
\addcontentsline{toc}{subsection}{可独立完成的完整试卷}


\begin{quote}
\footnotesize
\begin{verbatim}
07-期末复习PPT全覆盖综合卷02
  150 分钟，100 分；独立覆盖期末复习 PPT 一到七。

08-补充模拟卷04-运行时环境补缺
  120 分钟，100 分；独立覆盖运行时环境、栈帧、堆管理和垃圾回收。

09-补充模拟卷05-代码生成与寄存器分配补缺
  120 分钟，100 分；独立覆盖代码生成、基本块、DAG、描述符和寄存器分配。

10-补充模拟卷06-机器无关优化补缺
  120 分钟，100 分；独立覆盖数据流分析、支配关系和自然循环。

11-Slides补充模拟卷07-前端理论与PDA补缺
  120 分钟，100 分；独立覆盖 Slides 前端额外考点。

12-Slides补充模拟卷08-指令调度与支配边界补缺
  120 分钟，100 分；独立覆盖调度、Dom/PostDom/DF 和 SSA。

13-Slides补充模拟卷09-符号执行SAT与Tseitin补缺
  120 分钟，100 分；独立覆盖符号执行与 SAT。

14-Slides补充模拟卷10-过程间分析与指针分析补缺
  120 分钟，100 分；独立覆盖过程间分析与 Andersen 指针分析。

上述八套均包含：完整卷面说明、全部题目数据、100 分题目、单独分页的完整解析。
做题时不需要打开期末复习 PPT、Slides 或其他试卷。
\end{verbatim}
\end{quote}

\subsection*{专项补缺与背诵资料}
\addcontentsline{toc}{subsection}{专项补缺与背诵资料}


\begin{quote}
\footnotesize
\begin{verbatim}
第一组：期末复习 PPT 全覆盖
01-最简综合模拟卷
02-补充模拟卷01-自顶向下分析与LR概念补缺
03-补充模拟卷02-语义分析AST符号表补缺
04-补充模拟卷03-中间代码表示补缺
05-剩余知识点题型适配模拟卷
06-期末复习PPT名词解释
07-期末复习PPT全覆盖综合卷02
08-补充模拟卷04-运行时环境补缺
09-补充模拟卷05-代码生成与寄存器分配补缺
10-补充模拟卷06-机器无关优化补缺

第二组：Slides 额外考点补充
11-Slides补充模拟卷07-前端理论与PDA补缺
12-Slides补充模拟卷08-指令调度与支配边界补缺
13-Slides补充模拟卷09-符号执行SAT与Tseitin补缺
14-Slides补充模拟卷10-过程间分析与指针分析补缺
\end{verbatim}
\end{quote}

复习时先做第一组，确保期末复习 PPT 一到七全部过一遍；再做第二组，把 Slides 中没有被期末复习 PPT 系统覆盖的内容补齐。

\section*{二、期末复习 PPT 覆盖矩阵}
\addcontentsline{toc}{section}{二、期末复习 PPT 覆盖矩阵}


\begin{quote}
\footnotesize
\begin{verbatim}
期末复习一：词法解析与 LL(1)
覆盖文件：01、02、05、06、07
覆盖点：词素、词法单元、正则表达式、NFA、DFA、epsilon-closure、子集构造、DFA 最小化、递归下降、回溯、左递归、左公因子、FIRST、FOLLOW、预测分析表、LL(1) 判断。

期末复习二：自底向上文法解析与 LR
覆盖文件：01、02、05、06、07
覆盖点：移入-归约、句柄、最右推导逆过程、LR(0) 项、closure、GOTO、规范 LR(0) 项集族、ACTION/GOTO 表、Shift/Reduce 冲突、Reduce/Reduce 冲突、SLR。

期末复习三：语义分析
覆盖文件：01、03、05、06、07
覆盖点：SDD、SDT、属性、综合属性、继承属性、S 属性定义、L 属性定义、依赖图、拓扑求值、AST、类型表达式、符号表、受控副作用。

期末复习四：中间代码表示
覆盖文件：01、04、05、06、07
覆盖点：IR 层次、三地址码、四元式、三元式、间接三元式、AST 到 TAC、SSA、phi 函数、phi 消除、布尔表达式短路求值、回填、if/while/break/continue 翻译。

期末复习五：运行时环境
覆盖文件：06、07、08
覆盖点：静态区、栈区、堆区、活动树、活动记录、栈帧布局、SP、BP、参数传递、返回值、调用/返回代码序列、堆管理、碎片、Best-Fit、First-Fit、垃圾回收、可达性、引用计数、标记-清扫、压缩、拷贝、分代回收。

期末复习六：代码生成与寄存器分配
覆盖文件：06、07、09
覆盖点：指令选择、目标语言、指令代价、目标代码地址、基本块、流图、活跃性、下一次使用、DAG 基本块优化、寄存器描述符、地址描述符、getReg、简单代码生成、寄存器分配。

期末复习七：机器无关优化
覆盖文件：06、07、10
覆盖点：公共子表达式消除、复制传播、死代码消除、常量折叠、数据流分析框架、到达定值、活跃变量、可用表达式、gen/kill、IN/OUT、支配、后支配、直接支配、支配树、回边、自然循环、可归约流图。
\end{verbatim}
\end{quote}

\section*{三、Slides 额外考点覆盖矩阵}
\addcontentsline{toc}{section}{三、Slides 额外考点覆盖矩阵}


\begin{quote}
\footnotesize
\begin{verbatim}
Ch1 Introduction to Compilers
覆盖文件：11
覆盖点：编译器整体流程、前端/中端/后端分工、解释器与编译器对比、pass 与 phase 的区别、IR 在编译管线中的作用。

Ch3-1 Finite Automata / Ch3-2 Lexical Analysis
覆盖文件：01、07、11
覆盖点：DFA/NFA 形式化定义、正则表达式与自动机等价、DFA 最小化、DFA 等价性、词法分析器生成流程。

Ch4-1 Context-Free Language and Push-down Automata
覆盖文件：11
覆盖点：CFG、CFL、PDA、NPDA、二义性、CFL 闭包性质、正则语言与上下文无关语言的区别。

Ch4-2 Syntax Analysis
覆盖文件：01、02、07、11
覆盖点：消除左递归、预测分析、FIRST/FOLLOW、LL(1)、非递归预测分析器、PDA 与预测分析表关系、LR 基础。

Ch6-1 Intermediate Code Generation / Ch6-2 Translation into SSA
覆盖文件：04、07、12
覆盖点：IR、TAC、控制流翻译、SSA、phi、支配关系、Dominance Frontier、Iterated Dominance Frontier、phi 插入位置。

Ch8-1 Target Code Generation / Ch8-2 Translation out of SSA / Ch8-3 Register Allocation
覆盖文件：07、09、12
覆盖点：目标代码生成、phi 消除、寄存器分配、寄存器重命名与假依赖消除、后端综合题型。

Ch9-1 Dataflow Analysis
覆盖文件：10、12
覆盖点：数据流分析框架、支配相关概念与 SSA 之间的关系。

Ch9-2 Symbolic Execution
覆盖文件：13
覆盖点：符号执行、路径条件、路径敏感、流敏感、SAT、SMT、CNF、Tseitin Transformation、Equisatisfiability。

Ch10 Instruction Scheduling
覆盖文件：12
覆盖点：指令调度目标、True Dependence、Antidependence、Output Dependence、RAW、WAR、WAW、依赖图、寄存器重命名、基本块调度。

Ch12-1 Interprocedure Analysis
覆盖文件：14
覆盖点：过程间分析、调用图、上下文敏感、调用串、函数摘要、CFL reachability、匹配调用/返回。

Ch12-2 Pointer Analysis
覆盖文件：14
覆盖点：指针别名、may/must alias、points-to set、flow-sensitive、flow-insensitive、context-sensitive、Andersen 约束、约束图闭包。
\end{verbatim}
\end{quote}

\section*{四、2025 卷子对照}
\addcontentsline{toc}{section}{四、2025 卷子对照}


\begin{quote}
\footnotesize
\begin{verbatim}
2025 题 1：词法分析
覆盖文件：01、07、11

2025 题 2：语法分析 LL(1)
覆盖文件：01、02、07、11

2025 题 3：指令调度
覆盖文件：12

2025 题 4：支配问题、Dom Frontier
覆盖文件：10、12

2025 题 5：符号执行、CNF、Tseitin、Equisatisfiability
覆盖文件：13

2025 题 6：Andersen 指针别名
覆盖文件：14
\end{verbatim}
\end{quote}

\section*{五、最终覆盖结论}
\addcontentsline{toc}{section}{五、最终覆盖结论}


\begin{quote}
\footnotesize
\begin{verbatim}
01-10 合起来覆盖期末复习 PPT 一到七。
11-14 补齐 Slides 中未被期末复习 PPT 系统覆盖的考点。
01-14 合起来覆盖当前 Exam/2025.pdf 的六类题目，并覆盖 Slides 主目录与 handout 中的核心考点。
\end{verbatim}
\end{quote}
\end{document}
